\section{Research provenance and corrections to the exploratory model}
The input corpus comprises ten DOCX notebooks and one spreadsheet. The sequence begins with monetary finance plus MVI, adds tariffs and demurrage, and subsequently introduces an investment economy and usownership. The later discussion specifies sovereign money, gradual firm-level ownership and heterogeneous participant contributions. The accompanying audit identifies every input by filename and SHA-256 hash. The notebooks are not redistributed in the public-ready repository.

Several corrections are prerequisites rather than alternative preferences. Money growth is a percentage of a stock; government expenditure is a flow relative to GDP. A 10\% increase in a money stock does not automatically yield revenue equal to 10\% of GDP. Likewise, the national-accounts category ``taxes on production and imports'' includes domestic taxation, notably VAT; it cannot be treated as customs duties surviving the abolition of domestic taxes \citep{esamethod}. Tax receipts may replace issuance, but subtracting tax/GDP directly from inflation is dimensionally and economically incorrect.

The original investment projections also combine household or corporate investment income with government receipts. That would require a specified public ownership or tax channel. Purchasing an existing asset transfers money to a seller; it does not itself remove that money from the economy. The proposed administrative saving of 3--5\% of GDP depends on an assumed 10\%-of-GDP administrative base and assumed efficiency gains. It is not an observed saving, and eliminating public healthcare or pensions cannot be described simply as removing bureaucracy. The model below therefore retains a public-consumption block and does not claim that a universal cash transfer replaces all public services.

The workbook contains three stock arithmetic exercises. In the 41\% sheet, issuance of 41\% and destruction of 33\% of opening money imply 8\% net money growth. Calling this inflation additionally assumes unchanged real output and velocity. In the 12.21\% sheet, issuance less 6.25\% destruction implies 5.96\% net growth. At 2\% real growth and fixed velocity, exact inflation is $1.0596/1.02-1=3.882\%$. In the final sheet, the assumed 4\% real growth and 2\% inflation determine a burn rate by subtraction. This is a controller calculation, not a predicted inflation outcome. Exact nominal growth is 6.08\%, not 6\%, so the corresponding starting-stock burn with 12\% issuance is 5.92\%. All 33 populated stock-flow rows reconcile arithmetically; the issue is the economic interpretation.

Our stable hybrid explicitly relaxes the original aspiration to eliminate conventional taxation. A no-tax scenario is tested separately and fails under the specified closure. Equal citizen equity is also an external comparator, not a silent replacement of firm-specific usownership. These distinctions preserve the proposal's framing while making alternatives and corrections visible.

